\subsection*{MT-S508. One quadratic Hecke input and the exact automorphic central-order gate}

Retain the Calegari family of MT-S501: $L/\mathbb Q$ has group
$S_5$, complex conjugation has type $(12)(34)$, and the extension
is unramified at $5$ with Frobenius of the same type.  Let
$E=L^{A_5}$, let $F=L^{F_{20}}$ have degree $6$, and put $H=EF$.
Write $\eta$ for the quadratic character of $E/\mathbb Q$ and
$\eta_F$ for its restriction to $\operatorname{Gal}(L/F)$.
Then $H/F$ is quadratic, and its nontrivial Hecke character is
$\eta_F$.  Define the \emph{finite} $L$-functions
\[
 C(s)=L_{\mathbb Q}(s,\eta),\qquad
 B(s)=L_F(s,\eta_F)=\frac{\zeta_H(s)}{\zeta_F(s)}.
 \tag{S508.1}
\]
Both $B$ and $C$ are entire.  The all-prime induction identity
of MT-S506 gives
\[
             B(s)=C(s)L_W(s),\qquad
             \zeta_F(s)=\zeta_{\mathbb Q}(s)L_{W\eta}(s).
 \tag{S508.2}
\]
Let $\varpi$ be the particular cuspidal representation of
$\mathrm{GL}_5/\mathbb Q$ supplied by Calegari's construction,
and put $A(s)=L_f(s,\varpi)$.  This is its standard finite
$L$-function, which is entire.  Fix this comparison representation;
the statements below concern its exact transfer, not an arbitrary
choice of an automorphic representation of the same degree.
Here matching at every place means agreement of the standard
local $L$-factors, including the archimedean factor, in the
precise sense of the cited comparison theorem.

The only additional zero-location input needed here is
\[
  (\mathcal R_B)\qquad
  B(\sigma)\ne0\quad
  \bigl(0<\sigma<1,\ \sigma\ne\tfrac12,\ \sigma\in\mathbb R\bigr).
 \tag{S508.3}
\]
It concerns one quadratic Hecke character over one degree-six
field and only its real zeros.  GRH for this single Hecke
$L$-function implies $(\mathcal R_B)$; universal $M_H$ is a
sufficient input but is not required.  No Dirichlet GRH assumption
is needed for the central-order and automorphy assertions below.

Put
\[
 \beta=\operatorname{ord}_{1/2} B,\qquad
 b=\operatorname{ord}_{1/2} C,\qquad
 n=\operatorname{ord}_{1/2} A.
 \tag{S508.4}
\]
These are nonnegative integers.  Under $(\mathcal R_B)$ one has
\[
             \beta\ge b+n,
       \qquad \operatorname{ord}_{1/2}L_W=\beta-b\ge n\ge0.
 \tag{S508.5}
\]
Moreover the following are equivalent:
\[
 \begin{gathered}
       \beta\le b+n,\\
       \beta=b+n,\\
       \text{Calegari's exceptional set }\Omega\text{ is finite},\\
       \text{the specified }\varpi\text{ matches }W
             \text{ at every place}.
 \end{gathered}
 \tag{S508.6}
\]
Whenever these conditions hold, $\varpi\otimes\eta$ also matches
$W\eta$ at every place.  Thus every irreducible Artin factor of
$L/\mathbb Q$ is automorphic and every nontrivial one is entire.
If the comparison has not closed, its remaining gate is the
strict positive discrepancy
\[
                        \beta-b-n\ge1.
 \tag{S508.7}
\]
The older sufficient inequality $\beta\le b$ is a special case
of (S508.6).  The new condition retains possible central zeros
of the known automorphic numerator.

\paragraph{The exact classical comparison and its normalization.}
The classical construction and local comparison are Calegari,
\emph{The Artin conjecture for some $S_5$-extensions},
author-hosted publisher PDF: Theorem~1.1 p.~2, construction of
$\varpi$ p.~10, Theorem~4.7 p.~12, Lemma~4.8 pp.~13--14,
Landau Lemma~4.9 p.~14, Lemma~4.10 pp.~14--15, and the final
Theorem~1.2 proof p.~15.  This is the publisher numbering;
the last exact-comparison lemma is numbered 4.9 in the
arXiv v1 locators used in MT-S504.

We use the actual archimedean-completed quotient
\[
 R(s)=\frac{G_{\varpi}(s)A(s)}{G_W(s)L_W(s)},
 \tag{S508.8}
\]
where $G$ denotes the archimedean factor.  Including the usual
positive conductor powers would only multiply this by a
nowhere-zero entire function and would not change any order
or real-pole argument.  With the convention that a finite local
Euler factor is $P_p(p^{-s})^{-1}$, its correct factorization is
\[
 \begin{aligned}
 R(s)&=R_{\mathcal B}(s)
      \prod_{p\in\Omega}
             \left(\frac{1+p^{-s}}{1-p^{-s}}\right)^2,\\
 R_{\mathcal B}(s)&=
       \frac{G_{\varpi}(s)}{G_W(s)}
       \prod_{p\in\mathcal B}
             \frac{P_{W,p}(p^{-s})}{P_{\varpi,p}(p^{-s})}.
 \end{aligned}
 \tag{S508.9}
\]
Here $\mathcal B$ is the finite set of exceptional bad primes,
chosen disjoint from $\Omega$.  The source's displayed
$\mu_{\mathcal B}$ on publisher p.~13 reverses the two finite
polynomials relative to its subsequent reciprocal-Euler-factor
convention.  Equation (S508.9) fixes the direction directly from
the definition (S508.8); the source's holomorphic, nonzero
real-interval property is unaffected by this finite inversion.
In particular at $p\in\Omega$ its explicit polynomials are
\[
 P_{\varpi,p}(T)=(1+T^2)(1-T)^3,\qquad
 P_{W,p}(T)=(1+T^2)(1-T)(1+T)^2,
\]
which give exactly the displayed positive Euler quotient.

The cited comparison supplies two precise implications:
if $\Omega$ is infinite then $R$ has a pole at some real
$\sigma\in(0,1)$; if $\Omega$ is finite then the local factors
agree at all places, including the bad and archimedean ones.
The first uses its nonnegative Euler-series coefficients and
Landau's theorem; the second uses the cited finite-place
comparison with functional equations and twists.  The source's
final proof applies the same comparison to the quadratic twist.
We retain these classical results as inputs, not a new proof
of their automorphic construction.

The ratio $G_{\varpi}/G_W$ is holomorphic and nonzero on $(0,1)$.
For clarity, the Artin factor can be checked directly here.
Complex conjugation has trace $1$ on $W$, so its $+1,-1$
eigenspaces have dimensions $3,2$ and
\[
 G_W(s)=\Gamma_{\mathbb R}(s)^3\Gamma_{\mathbb R}(s+1)^2
       =\Gamma_{\mathbb R}(s)\Gamma_{\mathbb C}(s)^2.
 \tag{S508.10}
\]
Also $\eta$ is even, so $G_C(s)=\Gamma_{\mathbb R}(s)$.
The field $F$ has signature $(2,2)$, since the degree-six
permutation character $\mathbf1+W\eta$ takes value $2$ on
complex conjugation.  The real places split in $H/F$, because
$E$ is real.  Thus
$G_B(s)=\Gamma_{\mathbb R}(s)^2\Gamma_{\mathbb C}(s)^2
=G_C(s)G_W(s)$, consistently with (S508.2).
Here $\Gamma_{\mathbb R}(s)=\pi^{-s/2}\Gamma(s/2)$ and
$\Gamma_{\mathbb C}(s)=2(2\pi)^{-s}\Gamma(s)$.
All these factors are finite and nonzero at every real
$0<s<1$.  The source verifies that $G_{\varpi}$ likewise
uses standard real/complex gamma factors there.

\begin{proof}[Proof of the one-Hecke theorem]
First, the subgroup $F_{20}$ contains odd permutations, so
its intersection with $A_5$ has index two.  It follows that
$H=EF$ is the quadratic extension of $F$ cut out by the
nontrivial restriction $\eta_F$.  Artin induction, including
ramified factors, gives (S508.1)--(S508.2).

Substituting (S508.2) into the actual quotient yields
\[
          R(s)=U(s)\,\frac{A(s)C(s)}{B(s)},\qquad
          U(s)=\frac{G_{\varpi}(s)}{G_W(s)}.
 \tag{S508.11}
\]
For any real $0<\sigma<1$, the function $U$ is holomorphic
and nonzero.  Therefore
\[
 \operatorname{ord}_{\sigma}R
   =\operatorname{ord}_{\sigma}A
     +\operatorname{ord}_{\sigma}C
     -\operatorname{ord}_{\sigma}B.
 \tag{S508.12}
\]
This formula explicitly accounts for zeros of the automorphic
numerator; discarding them would give only the older sufficient
condition.

Under $(\mathcal R_B)$, the denominator in (S508.11) is nonzero
at every real point of $(0,1)$ other than the centre.
The other factors are holomorphic there, even if $C$ has a
zero.  Thus the only possible real pole of $R$ is at $1/2$,
where its order is $n+b-\beta$.
If $\Omega$ is infinite, the classical dichotomy forces a pole
there, hence $\beta>b+n$.  If $\Omega$ is finite, the exact
comparison gives $R\equiv1$ in the archimedean normalization
of (S508.8), hence $\beta=b+n$.  These two cases exhaust the
possibilities and prove (S508.5)--(S508.7).

The reverse implication in the last equivalence is also literal:
if all local factors of the specified comparison agree, then
no prime can have the two distinct polynomials displayed above,
so its exceptional set is empty.  Conversely, finite $\Omega$
forces exact local equality by the cited comparison.
For the quadratic twist, equality outside finitely many places
persists after twisting; the same exact-comparison argument,
as explicitly applied on source p.~15, identifies
$\varpi\otimes\eta$ with $W\eta$ at every place.  No second
Hecke zero-location hypothesis is needed for this automorphy
conclusion.  The dimensions $4$ and $6$ are already automorphic
in the specified family by Theorem~1.1, and the two linear
characters are classical.  The claimed holomorphy follows.
\end{proof}

\paragraph{The positive Euler quotient controls a whole half-plane.}
The same single real-zero hypothesis $(\mathcal R_B)$ implies
that $R$ is holomorphic and nonzero on $\Re s>1/2$.
In that half-plane it gives the exact divisor relation
\[
        \operatorname{ord}_z B
             =\operatorname{ord}_z A+\operatorname{ord}_z C.
 \tag{S508.13}
\]
In particular $L_W$ is holomorphic there and has the same
zero orders as $A$.  The poles of $L_W$ in the open critical
strip can only occur on its critical line, away from the real
axis.  This assertion requires neither $A_D$ nor full GRH for $B$.

\begin{proof}
Put
\[
 D_\Omega(s)=\prod_{p\in\Omega}
             \left(\frac{1+p^{-s}}{1-p^{-s}}\right)^2.
\]
This is the positive-coefficient Dirichlet series in the proof
of source Lemmas~4.8--4.9.  It is meromorphic on the plane
because $D_\Omega=R/R_{\mathcal B}$.  If $\Omega$ is infinite,
that proof identifies an abscissa of absolute convergence
$\gamma\in(0,1)$ and a pole at $\gamma$; the convergence at $1$
uses the stronger prime-sum bound proved in Theorem~4.7, not
Dirichlet density zero alone.  Since $R_{\mathcal B}$ is nonzero
and holomorphic at every real point of $(0,1)$, this is also a
pole of $R$.  Under $(\mathcal R_B)$, (S508.11) shows that this
pole can only be at $1/2$.  Hence $\gamma=1/2$.
If $\Omega$ is finite, the finite Euler product instead
converges for $\Re s>0$.

For every real $\sigma>1/2$ the nonnegative Dirichlet series
therefore converges, and its prime coefficients are $4$.
It follows that $\sum_{p\in\Omega}p^{-\sigma}<\infty$.
The logarithms of the Euler factors consequently converge
absolutely and locally uniformly on $\Re s>1/2$.
Each individual factor is nonzero in $\Re s>0$, so the
infinite product is holomorphic and nonzero on $\Re s>1/2$.
This also covers an empty or finite exceptional set.

The gamma quotient in $R_{\mathcal B}$ is holomorphic and
nonzero on all of $\Re s>0$: its standard gamma factors have
no poles or zeros there.  The finite polynomial factors are
also nonzero there, since their roots have modulus one and
$|p^{-s}|<1$.  Thus $R_{\mathcal B}$, and consequently $R$,
are holomorphic and nonzero on $\Re s>1/2$.
The same is true of $U$.  Taking orders in (S508.11) proves
(S508.13), and $L_W=UA/R$ proves the asserted holomorphy
and equality of zero orders.  The meromorphic Artin functional
equation for the selfdual $W$ reflects poles in
$0<\Re s<1/2$ to poles in $1/2<\Re s<1$;
its gamma and conductor factors are nonzero in this open strip.
Thus it excludes poles in the left half as well.  The centre
is already holomorphic by (S508.5).
\end{proof}

\paragraph{A single Hecke GRH also transfers to the known automorphic numerator.}
Assume now the full GRH for $B$, which implies
$(\mathcal R_B)$.  Then both $A=L_f(s,\varpi)$ and $C=L(s,\eta)$
satisfy GRH, whether or not the central equality (S508.6) holds.
Moreover the meromorphic divisor of $L_W=B/C$ in the open
critical strip is supported on the critical line.  This last
statement allows poles on that line; it does not prove Artin
holomorphy there.

\begin{proof}
For $1/2<\Re z<1$, GRH for $B$ says that
$\operatorname{ord}_z B=0$.  The two terms on the right of
(S508.13) are nonnegative, since $A$ and $C$ are holomorphic.
They are therefore both zero.  Their usual completed functional
equations then exclude zeros in $0<\Re z<1/2$ as well.
For the cuspidal unitary $\varpi$, the functional equation
pairs its standard $L$-function with that of its contragredient,
which is the complex conjugate under $s\mapsto\bar s$;
thus a zero is reflected across $\Re s=1/2$.
All gamma factors are nonzero in the open strip.  This proves
GRH for $A$ and $C$.  The factorization $L_W=B/C$ now proves
its critical-line divisor support.  If the central equality
also holds, the already proved exact transfer makes $L_W$
entire, so the divisor assertion is its ordinary GRH.
\end{proof}

This is specifically a \emph{one-$B$} transfer.  It is not a
claim that an $M_H$-conditional zero-location statement for
$\varpi$ was previously unavailable through its known base-change
product.  No GRH for $L_{W\eta},L_V,L_{V\eta}$ or $L_X$ follows
from the argument just given.  Once exact transfer is proved,
GRH for $\zeta_F$ separately gives GRH for $L_{W\eta}$ from
(S508.2), and universal $M_H$ is one sufficient way to supply
the other zero-location inputs.

\paragraph{The quadratic-twist expression for the same central gate.}
Put
\[
 A_\eta(s)=L_f(s,\varpi\otimes\eta),\qquad
 n_\eta=\operatorname{ord}_{1/2}A_\eta,
 \qquad f=\operatorname{ord}_{1/2}\zeta_F,
 \qquad k=\beta-b-n.
\]
Under $(\mathcal R_B)$ these nonnegative orders satisfy
\[
             k=n_\eta-f,\qquad 0\le k\le n_\eta.
 \tag{S508.14}
\]
Consequently $n_\eta\le f$ is equivalent to the same exact
transfer condition (S508.6).  In particular the additional
nonvanishing $A_\eta(1/2)\ne0$ suffices.  Neither that
nonvanishing nor this inequality is asserted unconditionally.

\begin{proof}
Every prime in $\Omega$ has Frobenius type $4\,1$, so
$\eta(p)=-1$.  Twisting its two displayed local polynomials
replaces $T$ by $-T$ and inverts their quotient.  Hence the
actual twisted completed ratio is
\[
 R_\eta(s)=\frac{G_{\varpi\eta}(s)A_\eta(s)}
                     {G_{W\eta}(s)L_{W\eta}(s)}
          =R_{\mathcal B,\eta}(s)D_\Omega(s)^{-1}.
 \tag{S508.15}
\]
The finite and archimedean factor $R_{\mathcal B,\eta}$ is
holomorphic and nonzero at every real $0<s<1$ before any
finiteness assertion about $\Omega$: $\eta$ restricts trivially
to $E$, so $(\varpi\otimes\eta)_E=\varpi_E$ has the same
known finite-image base change, and the local parameter argument
of source Lemma~4.8 still gives root-of-unity finite Euler
parameters after twisting.  In particular $\eta$ is even,
so it does not change the archimedean factors.  At the centre $R_{\mathcal B}$ and
$R_{\mathcal B,\eta}$ have order zero, and (S508.12) gives
$\operatorname{ord}_{1/2}D_\Omega=-k$.  Thus
$\operatorname{ord}_{1/2}R_\eta=k$.
On the other hand $\zeta_{\mathbb Q}(1/2)\ne0$, so (S508.2)
gives $\operatorname{ord}_{1/2}L_{W\eta}=f$.
Taking orders in the first quotient of (S508.15) gives
$k=n_\eta-f$.  Since $f\ge0$ and $k\ge0$ by (S508.5),
(S508.14) follows.  Equality $k=0$ is exactly (S508.6).
\end{proof}

\paragraph{Real-axis holomorphy needs only the one real-zero input.}
Under $(\mathcal R_B)$ every irreducible Artin factor in this
family is holomorphic at every real point of $(0,1)$.
For $W$, the half-plane argument and its functional equation
exclude every noncentral real pole, while (S508.5) handles
the centre.  For $W\eta$, (S508.2) and the elementary inequality
$\zeta_{\mathbb Q}(\sigma)<0$ for $0<\sigma<1$ prove
holomorphy on that real interval without any GRH input.
The other nontrivial factors are already entire by the
classical theorem.  Thus the earlier $A_D+(\mathcal R_B)$
corollary strengthens to $(\mathcal R_B)$ alone after the
positive Euler quotient is used.  This remains a conditional
statement about one Hecke function, not a verification of
$(\mathcal R_B)$.

\paragraph{The precise residual exceptional-prime law.}
If $k>0$, then $\Omega$ is infinite and, as $\epsilon\downarrow0$,
\[
 \sum_{p\in\Omega}p^{-1/2-\epsilon}
          =\frac{k}{4}\log\frac1\epsilon+O(1).
 \tag{S508.16}
\]
Thus an unclosed transfer forces an explicit positive logarithmic
singularity of the exceptional-prime harmonic sum, with the
same integer $k$ as the central excess.  This is an Abelian
statement at the Dirichlet-series boundary; no asymptotic
formula for the unweighted count of these primes is inferred.

\begin{proof}
By (S508.9), $D_\Omega$ has a pole of order $k$ at $1/2$.
It is positive for real $s>1/2$, so its leading Laurent
coefficient there is positive.  Consequently
\[
       \log D_\Omega(1/2+\epsilon)
             =k\log(1/\epsilon)+O(1).
\]
In the absolutely convergent half-plane the Euler logarithm is
\[
 \log D_\Omega(s)=4\sum_{p\in\Omega}p^{-s}+E(s),\qquad
 E(s)=4\sum_{p\in\Omega}\sum_{j\ge1}
              \frac{p^{-(2j+1)s}}{2j+1}.
 \tag{S508.17}
\]
The remainder is holomorphic for $\Re s>1/3$: on compact
subsets it is bounded termwise by a constant times
$\sum_p p^{-3\Re s}$, with locally uniform convergence.
In particular it is bounded near $s=1/2$.  Substituting the
real asymptotic for $\log D_\Omega$ into (S508.17) gives
(S508.16), with its exact coefficient $k/4$.
\end{proof}

\paragraph{Geometric and run attachment.}
The degree-six pushforward of the quadratic character of
$H/F$ is $\eta\oplus W$.  Equation (S508.12) isolates the
real divisor of the comparison between this Artin summand
and the known automorphic numerator.  The new gate is the
single central excess $\beta-b-n$, after the single real-zero
condition (S508.3); neither is evaluated by this proof.
This refines MT-S504/MT-S505 and the MT-S507 pole gate in
\texttt{GRH\_NIGHT/complexity}.  The source comparison remains
a classical input, $A_D$ remains the designated Dirichlet
input, and $M_H$ remains conditional.  No universal Artin
holomorphy, functoriality, or zero-independent positive
cohomology statement is inferred.
